\lysh{34354}{2}

\begin{alist}
    \item Ψάχνουμε το $f_3\%$. Γνωρίζουμε ότι ισχύει: $f_1\% + f_2\% + f_3\% + f_4\% + f_5\% = 100$ ή $10 + 10 + f_3\% + 35 + 15 = 100$ ή $70 + f_3\% = 100$ ή $f_3\% = 100 - 70 = 30$. Άρα $f_3\% = 30$.
    
    \begin{center}
    \begin{mytblr}{}
        Κλάσεις $[-)$ & Σχετική Συχνότητα $f_i\%$ \\        
        0-4 & 10 \\
        4-8 & 10 \\
        8-12 & \textbf{30} \\
        12-16 & 35 \\
        16-20 & 15 \\
    \end{mytblr}
    \end{center}

    \item Το ποσοστό των μαθητών που έγραψαν βαθμό μεγαλύτερο ή ίσο του 12 είναι:
    $35\% + 15\% = 50\%$, είναι δηλαδή το $50\%$ του συνόλου και επειδή είναι 42400 μαθητές άρα το σύνολο των μαθητών (το $100\%$) θα είναι: $42400 \cdot 2 = 84800$ μαθητές.

    \item Το ποσοστό των μαθητών που έγραψαν βαθμό μεγαλύτερο ή ίσο του 16 είναι:
    $15\%$. Το σύνολο των μαθητών είναι 84800 από ερώτημα β), συνεπώς το πλήθος των μαθητών που αναζητούμε είναι το $15\%$ του συνόλου, δηλαδή $15\% \cdot 84800 = 12720$ μαθητές.
\end{alist}
\vspace{6mm}
\noindent\rule{\textwidth}{0.4pt}
\vspace{6mm}

